\chapter{Category theory}

Category is the theory of structure. More generally, the form of generalization for mathematical structure. Such to say it borrows many formalism from different characteristics of mathematics, including the function descriptions, the object and class description, and so on. 

Such can be made into examples, of which typically can be sequences, of which is a finite list of elements of a certain type, denoted $[a_{0},\dots,a_{n-1}]$. The set of sequences over $A$ is denoted $Seq(A)$. Operation defined on them can be:
\begin{enumerate} 
    \item tips $=$ $a\mapsto [a]:A\to Seq(A)$. 
    \item cons (Notationally $\_ :\_$) $= (a,[a_{0},a_{1},\dots,a_{n-1}])\mapsto[a,a_{0},\dots,a_{n-1}]$. Formally, $A\times Seq(A)\to Seq(A)$.     
    \item joins $= ([a_{0},\dots,a_{m-1}],[a_{m},\dots,a_{n-1}])\mapsto[a_{0},\dots,a_{n-1}]$. Formally, $Seq(A)\times Seq(A)\to Seq(A)$. 
    \item $Seq(f)$ (function mapping) $= [a_{0},\dots,a_{n-1}]\mapsto[f(a_{0}),\dots,f(a_{n-1})]$. Formally, $Seq(A)\to Seq(B)$ whenever $f:A\to B$ is a morphism. 
    \item $\bigoplus /$ (direct sum) $=[a_{0},\dots,a_{n-1}]\mapsto a_{0}\oplus\dots \oplus a_{n-1}$. Formally, $Seq(A)\to A$ whenever $\bigoplus:A\times A\to A$. and $\bigoplus$ is associative and has a neutral element. 
\end{enumerate}
Wonder of what it is for the name, those examples are certain of the object behind the theory and form the normal viewport of category theory. 
\section{Introduction to Functions}

By the name, generally \textit{category theory} can be said as the generalization of mathematical structures, and their morphism, or functions in specific case, between elements of underlying mathematical objects, and between structures themselves. 

The first distinctive characteristic of category theory comes from the fact it lies solely on the \textit{abstractness} that it concurs with. For example, is the methodology of generalizing the relation and action between objects into two kinds - the operation and the function, or in to \textit{morphism}. 

A function, or mapping, is a kind of relation between elements in one 'object' to 'another object' of similar or different type. Often, in elementary mathematics, it would be different type of sets. For example, in the context of analysis, it would be $\mathbb{R}\longrightarrow \mathbb{R}$, or with $\mathbb{C}$. Given two objects, $A$ and $B$, the mapping from $A\longrightarrow B$ is denoted as $f: A\longrightarrow B$. The operation is unary (with single operand, and single output), thus we call function a unary operation on objects. To specify the input and output, we have the convention: 
\begin{equation*}
    A = src(f) = dom(f) \quad B = tgt(f) = range(f)
\end{equation*}

A successive sequence of function, for example, between three objects $(A,B,C)$ expressed as 
\begin{equation*}
    A \underset{f}{\longrightarrow} B \underset{f}{\longrightarrow} C 
\end{equation*}
can have a composite function, $g \circ f$ (here we have $g$ before $f$, since in the commutative diagram, $range(f)=dom(g)$, hence $g$ 'wraps' above $f$):

\begin{center}
    \begin{tikzcd}
        A \arrow[r, "f"] \arrow[rd, "g \circ f"', dotted] & B \arrow[d, "g"] \\
                                                          & C               
        \end{tikzcd}
\end{center}

This function, $g \circ f: A \longrightarrow B$ is written as
\begin{equation*}
    g \circ f (a) = g(f(a))
\end{equation*}
Further operation chain including the quadruplet $(A,B,C,D)$, with $h:C \longrightarrow D$ with such 
\begin{center}
    \adjustbox{scale=1}
    {
        \begin{tikzcd}
            A \arrow[r, "f"] \arrow[rd, "g \circ f"'] & B \arrow[d, "g", no head, Rightarrow] \arrow[rd, "h \circ g", dashed] &   \\
                                                      & C \arrow[r, "h"']                                                     & D
            \end{tikzcd}
    }
\end{center}

Here, the formation of $h\circ g$ and $g \circ f$ let us compare $(h\circ g)\circ f$ and $h \circ (g \circ f)$ as indicated in the diagram. The result is quite surprising, $(h\circ g)\circ f = h \circ (g \circ f)$. This is because accordingly, 
\begin{equation*}
    ((h\circ g) \circ f)(a) = h(g(f(a))) = (h\circ (g \circ f))(a)
\end{equation*}

For any two functions to be "equal", for every argument, their value must be the same. Under the observation of commutative diagram as above, it is analogous to having the same endpoint. The observation above is also applied for identity function, $1_A : A \longrightarrow A$ given by $1_A (a)=a$. Theoretically, those function acts as units for the composition operation: $f \circ 1_A = f = 1_A \circ f$. The commutative diagram for identity of the tuple $(A,B)$ is

\begin{center}
    \adjustbox{scale=1}
    {
        \begin{tikzcd}
            A \arrow[r, "1_A"] \arrow[rd, "f\circ 1_A"] & A \arrow[rd, "1_B \circ f"] \arrow[d, "f" description] &   \\
                                                        & B \arrow[r, "1_B"]                                     & B
            \end{tikzcd}
    }
\end{center}

\section{Category and Precategory}
We now give a formal definition of a category. 
\begin{definition}
    A category is a system of: 
    \begin{description}
        \item[Objects] denoted by $A,B,C,\dots$.
        \item[Morphism] denoted by $f,g,h,\dots$.
        \item[Relation on Morphism/Objects] called typing of the morphisms. We say that for $f:A\to B$, then $A\to B$ is the type of $f$, and $f$ is a morphism from $A$ to $B$. Then $\mathrm{src}(f)=A$ and $\mathrm{tgt}(f)=B$.
        \item[Composition] A binary partial operation denoted as $f \circ g=gf=f;g$.
        \item[Identity] For each object $A$ a distinguished morphism, called \textbf{identity} on $A$. By default, we denote it as $id_{A}$, denotes the identity on object $A$.  
    \end{description}
\end{definition}



This category definition defines the categorical language in which properties of the category can be stated, of which can be built from normal logical connectives and quantification and equality. Of course, category does have axioms, as for all formalism of mathematical system, needs axioms for its description of the mathematical structure. 

\begin{theorem}[Axiom of Category]
For a category, there are 5 axioms of which from (1) to (3) are of typing, and the others are for morphisms.\footnote{A morphism term $f$ is \emph{well-typed} when the typing can be derived for some objects $A,B$ according to these axioms.}

\begin{enumerate}
    \item[(I)] If $f:A \to B$ and $f:A' \to B'$, then $A = A'$ and $B = B'$ \hfill (unique-Typing).
    \item[(II)] If $f:A \to B$ and $g:B \to C$, then $f \circ g : A \to C$ \hfill (composition-Typing).
    \item[(III)] $id_A : A \to A$ \hfill (identity-Typing).
    \item[(IV)] $(f \circ g) \circ h = f \circ (g \circ h)$ \hfill (composition-Association).
    \item[(V)] $id \circ f = f = f \circ id$ \hfill (Identity).
\end{enumerate}
\end{theorem}

A \textbf{precategory}\index{precategory} then is a category of which the $(I)$ requirement is dropped. By simple trick (as is stated by the original author), we can construct a category $\mathcal{B}$ from a pre-category $\mathcal{A}$. Suppose that $\mathcal{B}$ is defined by: 
\begin{itemize}
    \item An object in $\mathcal{B}$ is an object in $\mathcal{A}$. 
    \item A morphism in $\mathcal{B}$ is a triple $(A,f,B)$ with $f:A\longrightarrow_{\mathcal{A}}B$. 
    \item $f:A\longrightarrow_{\mathcal{B}}B\equiv A=A',B=B'$ where $f=(A',f',B')$. $(*)$
    \item $f\circ_{\mathcal{B}} g=(A,f'\circ_{\mathcal{A}}g',C)$ where $(A,f',B)=f$ and $(B,g',C)=g$. 
    \item $id_{\mathcal{B},A}=(A,id_{\mathcal{A},A},A)$. 
\end{itemize}
Then, we can deduce or prove that $\mathcal{B}$ is a category. Our best concern of this proof will be the starred $(*)$ statement of which validate the notion of such. Using the notation, we see that a morphism $f=(A,f',B)$ of which $f':A'\to_{\mathcal{A}} B'$. Then, of which we can see that for $f_{\mathcal{B}}=f'_{\mathcal{A}}$, then the morphism $f:A\to_{\mathcal{B}}B$ of which is defined on $\mathcal{B}$, will be also given of such definition, that is: $f:A\to_{\mathcal{B}}B=(A,f',B)$, but $f_{\mathcal{B}}=f'_{\mathcal{A}}$, then both object of which also, again, $A,B\in \mathcal{B}\land A,B\in \mathcal{A}=T$ (both of them is in $\mathcal{A}$ and $\mathcal{B}$), then it must be said that $\mathcal{B}$ is a category since this satisfy the axiom $(I)$. 

The axioms on the morphisms and composition are very weak, so that many mathematical structures can be rendered as a category. By imposing extra axioms, still in the categorical language, the categories might have more of the properties of interest. A  Cartesian closed category  (CCC) is a category in which the extra properties make the morphisms behave like real functions, that is, there exists the notion of currying and of applying a curried morphism. In brief, currying is a process of transforming an operation on two variables into an operation on one variable that returns a function taking the second variable as an argument. For example, for a function $f:X\times Y\to Z$, currying will create a function $\hat{f}:X\to Z^{Y}$ of which $\hat{f}(x)=(y\mapsto f(x,y))$. One can also think about this as being $X\times Y\to Z\equiv X\to[Y\to Z]$. There is a close relationship between this type of category and typed $\lambda$-calculi. 

Further into notation, given two objects $A$ and $B$ in an arbitrary category $\mathcal{C}$, then the collection of morphisms is denoted $\mathcal{C}[A,B]$ or $\mathcal{C}(A,B)$. This is called a hom-set, and another frequent notation is $\mathrm{Hom}_{\mathcal{C}}[A,B]$.  

\section{Functors}

Further increment in the scope of our problems, we are interested in morphisms between categories. A homomorphism of categories is called a functor.

\begin{definition}
    A functor 
    \begin{equation*}
        F: \mathbf{C} \longrightarrow \mathbf{D}
    \end{equation*}
    between category $\mathbf{C}$ and $\mathbf{D}$ is a mapping of objects to objects and morphisms to morphisms such that 
    \begin{enumerate}
        \item $F(f:A\longrightarrow B)=F(f):F(A)\longrightarrow F(B)$
        \item $F(1_A) = 1_{F(A)}$
        \item $F(g \circ f)=F(g)\circ F(f)$
    \end{enumerate}
\end{definition}

The formula $F: A\longrightarrow B$ means that $F$ is a functor 

Retrospectively, $F$ preserves domains and codomains, identity morphism, and binary composition. The functor $F:\mathbf{C}\longrightarrow \mathbf{D}$ thus gives a sort of "picture" of $\mathbf{C}$ and $\mathbf{D}$. Commutative diagram is shown as below: 

\begin{center}
    \adjustbox{scale=1}{
    \begin{tikzcd}
        \mathbf{C} \arrow[ddd, "F"] & A \arrow[rd, "g \circ f"] \arrow[r, "f"]          & B \arrow[d, "g"]       \\
                                    &                                                   & C                      \\
                                    &                                                   & F(B) \arrow[d, "F(g)"] \\
        \mathbf{D}                  & F(A) \arrow[ru, "F(f)"] \arrow[r, "F(g \circ f)"] & F(C)                  
        \end{tikzcd}
}
\end{center}


One can clearly see that each category have the identity morphism $1_\mathbf{C}:\mathbf{C}\longrightarrow \mathbf{C}$. Regarding this, considering only two categories, we have 
\begin{center}
    \adjustbox{scale=1.1}{
        \begin{tikzcd}
            \mathbf{C} \arrow[rr, "F"] \arrow["1_{\mathbf{C}}"', loop, distance=2em, in=305, out=235] &  & \mathbf{D} \arrow[loop, distance=2em, in=305, out=235] \arrow[loop, distance=2em, in=305, out=235] \arrow[loop, distance=2em, in=305, out=235] \arrow["1_{\mathbf{D}}"', loop, distance=2em, in=305, out=235]
            \end{tikzcd}
    }
\end{center}

Hence, the components for a category between categories is sufficed. We call this $\mathbf{Cat}$ category to represent the category of all categories and functors. 

\section{Examples of Category}

There are many examples of categories, which we will discuss as a practice of defining mathematical system in categorical language. 

\subsection{Unital Ring}

For a unital ring $R$, the category $\operatorname{Mat}_R$ is category whose objects are positive integers and in which the set of morphisms from $n \to m$ is the set of $m \times n$ matrices with value in $R$. Composition is by matrix multiplication
\begin{center}
    \adjustbox{scale=1.1}
    {
        \begin{tikzcd}
            n \arrow[rd, "A"] \arrow[rr, "A \cdot B"] \arrow["1_n"', loop, distance=2em, in=215, out=145] &                                                                       & k \arrow["1_k"', loop, distance=2em, in=35, out=325] \\
                                                                                                          & m \arrow[ru, "B"] \arrow["1_m"', loop, distance=2em, in=305, out=235] &                                                     
            \end{tikzcd}
    }
\end{center}

\subsection{Group and Monoid}
A group or generally, a monoid defines the category $BG$ with a single object. The group elements are its morphisms, each group elements represent a distinct endomorphism (as we will discuss later) with composition given by multiplication. The identity element $e\in G$ acts the identity morphism for the unique object of the category. 

Generally, monoid is a set $M$ equipped with a binary operation $M \times M \longrightarrow M$ and a distinguished element $u\in M$ such that: 
\begin{equation*}
    \forall (x,y,z)\in M\quad \begin{cases}
        u \cdot x = x = x \cdot u \\
        x \cdot ( y \cdot z) = (x \cdot y) \cdot z
    \end{cases}
\end{equation*}

Equivalently, a monoid is a category with just one object. The arrow of the category are the elements of the monoid. In particular, the identity arrow is the unit element $u$. 

\section{Investigation on Morphism}

In modern mathematics, we have different way of formalising the mathematical foundation. While listing the object is interesting, all the axiom focus on sorting objects fitting to different categories, based on their behaviour within such system. For example, for the monoid, it is a tuple $(M,e,\times)$ that satisfies the above notion. 

Similar, our main investigation would be mostly morphism and functors, since from investigating their behaviour, can we achieve understanding the underlying abstraction of object. 

\subsection{Total Morphism}

A total morphism is a general morphism, or a bare morphism. It is similar to the aforemention morphism, which takes the form $M_{\mathcal{B}}:(O_{1},\dots,O_{n})\longrightarrow (O_{1}',\dots O_{n}'$. Usually, it takes the form of a mapping size $n\to m$. In the context of linear algebra, such morphism can be said as a linear mapping of $M_{\mathbf{Lin}}: (\mathbb{F}_1, \dots, \mathbb{F}_n)\longrightarrow (\mathbb{F}_{1}',\dots,\mathbb{F}_{m}')$

\subsection{Homomorphism}

In a restrictive sense, a homomorphism is a function between two algebras that preserves its algebraic structure. It makes sense for category of algebraic type to be called a structure set, like with the case of monoid and groups. Then, a homomorphism is a synonym for morphism (remember that total morphism takes an arbitrary domain space and returns also an arbitrary domain space) when structure sets are generalized to arbitrary object. 

Traditionally, a homomorphism between two magmas (structured set equipped with binary operations on itself) $A$ and $B$ is a function 
\begin{equation*}
    \phi : A \longrightarrow B
\end{equation*}

That satisfy the underlying binary operation. Specifically, for all $a_{1}, a_{2}\in A$, we have 
\begin{equation*}
    \phi (a_1 \times a_2) = \phi (a_1) \times \phi(a_2)
\end{equation*}

This definition gives us the correct notion of \textbf{magma, semigroup} and \textbf{group} homomorphism, but not for monoid homomorphism and ring homomorphism. 

\begin{definition}[monoid/ring homomorphism]
    A homomorphism between two monoids $A$ and $B$ is a semigroup homomorphism, $\theta: A \longrightarrow B$ of the underlying semigroups that preserves identity element, that is: 
    \begin{equation*}
        \theta(1_{A}) = 1_{B}
    \end{equation*}
    A ring homomorphism is a function between rings that is a homomorphism for both the additive group and the multiplicative monoid.   
\end{definition}

More generally, a homomorphism between sets equipped with any algebraic structure is a map preserving this structure. This can be made precise using Lawvere theories, but generally, structure-preserving can be made clear using intuitional explanation. 

There are many structures in category theory. Suppose with a triplet $(x,y,z)$ where, $x,y,z\in X$, of which on $X$ addition is defined. This means that some triple $(x,y,z)$ of the elements of $X$ are special so far as $x+y=z$ is true. Then we can say that the triplet is a structure that conforms onto elements of $X$, and we denote it as $(x,y,z)\in plus$. Continuing with the assumption for the second set $Y$ of which satisfy certain requirements as $X$, essentially the same construct, and with the second structure named $plus'$. 

From $X$ to $Y$, we define a morphism $\phi:X \longrightarrow Y$. We said that $\phi$ is a structure preserving morphism, or homomorphism, if 
\begin{equation*}
    (x_1,x_2,x_3) \in plus \underset{\phi}{\longrightarrow} (\phi(x_1),\phi(x_2),\phi(x_3)) \in plus'
\end{equation*}
\subsection{Isomorphism}
Isomorphism is a somewhat expansion from homomomorphism, as far as the definition is concerned of. 

By the standard group theory definition, we can say about isomorphism as followed.
\begin{definition}[group isomorphism]
    Let $G$ and $ H $ be groups. An isomorphism $ f: G \longrightarrow H $ is a bijection $ f:G \longrightarrow H $ such that for all $g_1, g_2 \in G$, 
    \begin{equation*}
        f(g_1 g_2) = f(g_1) f(g_2)
    \end{equation*}
\end{definition}
From the definition of group isomorphism, isomorphism can be thought as the bijective version of homomorphism, that is, as the previous intuitive explanation of homomorphism structure preservation property, then 
\begin{equation*}
    (x_1,x_2,x_3) \in plus \underset{\phi}{\iff} (\phi(x_1),\phi(x_2),\phi(x_3)) \in plus'
\end{equation*}
for an isomorphism $\phi: A \longrightarrow B$.

For a category-theoretic definition, we have the following: 

\begin{definition}[categorical isomorphism]
    In any category $\mathbf{C}$, an arrow $f:A\to B$ is called an isomorphism if there is an arrow $g: B \to A$ in $\mathbf{C}$ such that: 
    \begin{equation*}
        g \circ f = 1_A \quad \text{and} \quad f \circ g = 1_B
    \end{equation*}
\end{definition}

effectively, this means that the arrow $f$ has its inverse. Since inverse are unique, we write $g=f^{-1}$. 